% \iffalse meta-comment
%
% Copyright (C) 2026 by the banjofret authors
%
% This work may be distributed and/or modified under the conditions of
% the LaTeX Project Public License (LPPL), either version 1.3c of this
% license or (at your option) any later version.  The latest version of
% this license is in https://www.latex-project.org/lppl.txt and version
% 1.3c or later is part of all distributions of LaTeX version 2008-05-04
% or later.
%
% This work has the LPPL maintenance status `maintained'.
%
% \fi
%
% \iffalse
%<*driver>
\ProvidesFile{banjofret.dtx}[2026/09/28 v1.0.0 banjofret documentation driver]
\documentclass[a4paper]{ltxdoc}
\IfFileExists{lmodern.sty}{\usepackage[T1]{fontenc}\usepackage{lmodern}}{}
\usepackage[a4paper,margin=2.3cm,marginparwidth=2cm]{geometry}
\usepackage{banjofret}
\usepackage{booktabs}
\usepackage{pdflscape}
\usepackage{listings}
\usepackage[colorlinks,linkcolor=blue!50!black,urlcolor=blue!50!black]{hyperref}
\lstdefinestyle{bf}{language={[LaTeX]TeX},basicstyle=\ttfamily\small,
  frame=single,rulecolor=\color{black!25},backgroundcolor=\color{black!3},
  keywordstyle=\color{blue!55!black},commentstyle=\color{black!55},
  xleftmargin=0pt,framexleftmargin=4pt,aboveskip=6pt,belowskip=4pt,
  columns=fullflexible,keepspaces=true,showstringspaces=false}
% Show a snippet (source listing) followed by its live result.
\newcommand\BFexample[1]{%
  \IfFileExists{examples/#1.tex}{\def\BFpath{examples/#1.tex}}{\def\BFpath{#1.tex}}%
  \par\medskip\noindent
  \begingroup\lstinputlisting[style=bf]{\BFpath}\endgroup
  \MakePercentComment\input{\BFpath}\MakePercentIgnore
  \par\medskip}
\renewcommand\MacroFont{\fontencoding\encodingdefault\fontfamily\ttdefault\fontseries\mddefault\fontshape\updefault\footnotesize}
\DisableCrossrefs
\CodelineIndex
\begin{document}
  \DocInput{banjofret.dtx}
\end{document}
%</driver>
% \fi
%
% \CheckSum{0}
%
% \GetFileInfo{banjofret.dtx}
%
% \title{The \textsf{banjofret} package\\[2pt]
%   \large Colour-coded fretboard diagrams for the 5-string banjo}
% \author{banjofret contributors}
% \date{Version 1.0.0 \,---\, 2026/09/28}
% \maketitle
%
% \begin{abstract}
% \noindent
% \textsf{banjofret} typesets a clean, black-and-white diagram of a 5-string banjo
% fretboard (24 frets, short fifth string starting at fret~5) with the chromatic
% note name printed on every string/fret position.  You then \emph{paint} the
% notes you care about --- single positions such as \verb|(3,5)|, whole
% strings, whole frets, or every occurrence of a given note --- using a small
% palette of minimalist colours chosen to support learning.  Note labels are
% always drawn on top of the strings and of the colours.  It is built on
% Ti\emph{k}Z and expl3 and works with pdf\LaTeX, Xe\LaTeX{} and Lua\LaTeX.
% \end{abstract}
%
% \tableofcontents
%
% \section{Introduction}
%
% The goal of the package is to let you design \emph{your own} banjo diagrams ---
% scales, chords, arpeggios, licks, positions --- with as little typing as
% possible.  There is one template (the fretboard) and a handful of commands to
% paint it:
%
% \begin{itemize}
%   \item you always get the full chromatic fretboard, open-G tuned, in black and white;
%   \item you paint by position, by string, by fret or by note name;
%   \item colours are short names (\texttt{blue}, \texttt{root}, \texttt{amber}, \dots);
%   \item note names are drawn last, so they stay legible over strings and colours.
% \end{itemize}
%
% \subsection{Requirements}
%
% \LaTeX{} 2022-06-01 or newer (for the built-in \verb|\NewDocumentCommand|), plus
% the packages \textsf{tikz}, \textsf{xcolor} and \textsf{l3keys2e}, all part of
% every standard \TeX~Live and MiK\TeX{} installation.  The package loads
% \textsf{tikz} (with the \texttt{calc} library) and \textsf{xcolor} itself.
%
% \subsection{Installation}
%
% Run \texttt{latex banjofret.ins} to generate \texttt{banjofret.sty}, and place
% it (together with this PDF) in your local \texttt{texmf} tree, for instance
% \texttt{TEXMF/tex/latex/banjofret/}.  For a one-off project it is enough to
% keep \texttt{banjofret.sty} next to your \texttt{.tex} file.  On Overleaf,
% upload \texttt{banjofret.sty} to the project.
%
% \section{Quick start}
%
% \begin{verbatim}
% \documentclass{article}
% \usepackage[landscape,margin=1.5cm]{geometry}
% \usepackage{banjofret}
% \begin{document}
%
% \begin{banjofret}[title={My first diagram}]
%   \bfpaint{(3,5)}{blue}          % string 3, fret 5
%   \bfpaintnote{G, B, D}{root}    % every G, B and D on the neck
% \end{banjofret}
%
% \end{document}
% \end{verbatim}
%
% With no painting commands at all you get the empty template with all
% 24~frets.  It is wide, so landscape pages (or a 12-fret window, see
% Section~\ref{sec:options}) are recommended:
%
% \begin{landscape}
% \BFexample{snip-empty}
% \end{landscape}
%
% \section{How the diagram is organised}
% \label{sec:coords}
%
% \paragraph{Strings.}
% Strings are numbered as on a banjo: string~1 is the thinnest of the four
% full-length strings (high D) and string~4 the thickest (low D).  String~5
% is the short drone string.  By default string~1 is drawn on top and string~5
% at the bottom, like in banjo tablature (\texttt{string order=tab}); use
% \texttt{string order=reverse} for the opposite.
%
% \paragraph{Frets.}
% Fret~0 is the open string and is drawn in the column left of the nut (the
% thick vertical bar).  Frets~1 to~24 follow.  Frets have equal width: the diagram
% is a schematic and not a scale drawing.
%
% \paragraph{The fifth string.}
% The fifth string starts at fret~5, where its tuning peg sits.  The disc at
% position \verb|(5,5)| is the open fifth string (G); \verb|(5,6)| is G$\sharp$/A$\flat$, and so on.  Positions \verb|(5,0)|
% to \verb|(5,4)| do not exist and are never drawn.
%
% \paragraph{Position notation.}
% A position is written \verb|(string,fret)|.  Thus \verb|(3,5)| means string~3
% at fret~5.  Several positions are separated by \emph{semicolons}:
% \verb|(3,5); (2,1); (1,0)|.
%
% \paragraph{Default tuning.}
% The default is standard open-G (G4--D3--G3--B3--D4):
%
% \begin{center}
% \begin{tabular}{@{}cccccc@{}}
% \toprule
% String & 1 & 2 & 3 & 4 & 5 \\
% \midrule
% Note   & D & B & G & D & G \\
% \bottomrule
% \end{tabular}
% \end{center}
%
% Other tunings can be selected with the \texttt{tuning} option (see below).
%
% \section{Painting commands}
%
% All painting commands must be used \emph{inside} a \texttt{banjofret}
% environment.  If two commands paint the same position, the \emph{last} one wins,
% so you can first paint a broad set of notes and then override individual
% positions.  The last argument of every command is a colour
% (Section~\ref{sec:colours}).
%
% \DescribeEnv{banjofret}
% \begin{quote}
% \verb|\begin{banjofret}[|\meta{options}\verb|]| \meta{painting commands}
% \verb|\end{banjofret}|
% \end{quote}
% Typeset one fretboard diagram.  The optional argument takes the options of
% Section~\ref{sec:options}.
%
% \DescribeMacro{\bfpaint}
% \begin{quote}
% \verb|\bfpaint{|\meta{positions}\verb|}{|\meta{colour}\verb|}|
% \end{quote}
% Paint one or more \verb|(string,fret)| positions.  Parentheses are optional;
% positions are separated by semicolons.
%
% \BFexample{snip-paint}
%
% \DescribeMacro{\bfpaintnote}
% \begin{quote}
% \verb|\bfpaintnote{|\meta{notes}\verb|}{|\meta{colour}\verb|}|
% \end{quote}
% Paint \emph{every} occurrence of one or more notes on the neck.  Notes are
% comma-separated names: a letter \texttt{A}--\texttt{G} followed by an optional
% accidental, which can be written \texttt{s} or \verb|#| for sharp and
% \texttt{b} for flat: \texttt{C}, \texttt{Fs}, \verb|F#|, \texttt{Bb}.  Enharmonic
% spellings are equivalent (\texttt{Fs} and \texttt{Gb} paint the same discs).
% This is the fastest way to draw scales and arpeggios:
%
% \BFexample{snip-gscale}
%
% \DescribeMacro{\bfpaintstring}
% \begin{quote}
% \verb|\bfpaintstring{|\meta{string}\verb|}{|\meta{frets}\verb|}{|\meta{colour}\verb|}|
% \end{quote}
% Paint a run of frets on one string.  The fret range is either a single fret
% (\texttt{5}) or a range (\texttt{2-7}).
%
% \DescribeMacro{\bfpaintfret}
% \begin{quote}
% \verb|\bfpaintfret{|\meta{fret range}\verb|}{|\meta{colour}\verb|}|
% \end{quote}
% Paint the same fret, or a range of frets, on all strings.
%
% \BFexample{snip-ranges}
%
% \DescribeMacro{\bflegend}
% \begin{quote}
% \verb|\bflegend{|\meta{colour}\verb|=|\meta{label}\verb|, |\dots\verb|}|
% \end{quote}
% Add a legend line below the diagram.  Use braces around labels that contain
% commas or equal signs.
%
% \section{Chords, scales and arpeggios}
%
% Chord shapes are most readable as position lists with the chord-tone colours
% \texttt{root}, \texttt{third}, \texttt{fifth} and \texttt{seventh}:
%
% \BFexample{snip-cchord}
%
% Whole-neck views work best with \texttt{focus}, which greys out every note that
% is not painted, so the eye finds the target notes at once:
%
% \BFexample{snip-triad}
%
% \section{Colours}
% \label{sec:colours}
%
% The colour argument of every painting command can be one of the
% palette names below.  The names are resolved to the package colours
% \texttt{bf}\meta{name} (for instance \texttt{blue} $\to$ \texttt{bfblue}),
% which can therefore also be used with any standard \texttt{xcolor} command.
% Any other \textsf{xcolor} expression (\texttt{red!30}, a colour defined with
% \verb|\definecolor|, \dots) is accepted as it is.
%
% \begin{center}
% \renewcommand{\arraystretch}{1.25}
% \begin{tabular}{@{}l c l l@{}}
% \toprule
% Name & Swatch & HTML & Suggested meaning \\
% \midrule
% \texttt{blue}    & \fcolorbox{black}{bfblue}{\rule{1cm}{0pt}\rule{0pt}{1.6ex}}   & \texttt{7DB4E8} & calm, main information\\
% \texttt{green}   & \fcolorbox{black}{bfgreen}{\rule{1cm}{0pt}\rule{0pt}{1.6ex}}  & \texttt{86CFA3} & targets, resolution, ``safe'' notes\\
% \texttt{amber}   & \fcolorbox{black}{bfamber}{\rule{1cm}{0pt}\rule{0pt}{1.6ex}}  & \texttt{F5C563} & attention, passing emphasis\\
% \texttt{coral}   & \fcolorbox{black}{bfcoral}{\rule{1cm}{0pt}\rule{0pt}{1.6ex}}  & \texttt{F29E8E} & the root, the most important note\\
% \texttt{violet}  & \fcolorbox{black}{bfviolet}{\rule{1cm}{0pt}\rule{0pt}{1.6ex}} & \texttt{B9A3E3} & colour tones, extensions\\
% \texttt{teal}    & \fcolorbox{black}{bfteal}{\rule{1cm}{0pt}\rule{0pt}{1.6ex}}   & \texttt{7BCFCF} & second family of tones\\
% \texttt{gray}    & \fcolorbox{black}{bfgray}{\rule{1cm}{0pt}\rule{0pt}{1.6ex}}   & \texttt{C9CED6} & neutral, context, ghost notes\\
% \bottomrule
% \end{tabular}
% \end{center}
%
% \noindent Semantic aliases make diagrams self-consistent across a whole book:
%
% \begin{center}
% \begin{tabular}{@{}l l @{\qquad} l l@{}}
% \toprule
% Alias & = & Alias & = \\
% \midrule
% \texttt{root}      & \texttt{coral}  & \texttt{seventh}   & \texttt{violet} \\
% \texttt{third}     & \texttt{amber}  & \texttt{extension} & \texttt{teal}   \\
% \texttt{fifth}     & \texttt{blue}   & \texttt{passing}   & \texttt{gray}   \\
% \texttt{target}    & \texttt{green}  &                    &                 \\
% \bottomrule
% \end{tabular}
% \end{center}
%
% \paragraph{Design rationale.}
% The palette follows a few well-established guidelines from the literature on
% multimedia learning and colour accessibility (they are guidelines, not laws):
% (1)~a \emph{small} set of hues, so that the meaning of each colour can be
% remembered; (2)~one \emph{consistent} meaning per colour --- the warm coral is
% kept for the most salient item, the root, while cool colours carry supporting
% information; (3)~soft, mid-light tones with black text, which keeps contrast
% high and the page calm, and avoids the visual vibration of pure primaries;
% (4)~hues that stay distinguishable under the common forms of colour-vision
% deficiency; and (5)~redundant coding: the note \emph{name} is always printed
% inside the disc and a legend can be added, so colour is never the only channel.
%
% \paragraph{Your own colours.}
% Define them with \textsf{xcolor} and use them by name:
% \begin{quote}
% \verb|\definecolor{mine}{HTML}{FFB3C1}|\\
% \verb|\bfpaint{(3,5)}{mine}|
% \end{quote}
%
% \section{Options}
% \label{sec:options}
%
% Options are given in the optional argument of the environment.  To change a
% default for the whole document, pass them as package options, for example
% \verb|\usepackage[last fret=12]{banjofret}|, or set them in the preamble with:
%
% \DescribeMacro{\banjofretset}
% \begin{quote}
% \verb|\banjofretset{|\meta{options}\verb|}|
% \end{quote}
%
% \begin{center}
% \renewcommand{\arraystretch}{1.2}
% \begin{tabular}{@{}l l p{7.6cm}@{}}
% \toprule
% Option & Default & Meaning\\
% \midrule
% \texttt{first fret}   & \texttt{0} & First fret shown. \texttt{0} shows the open-string column and the nut.\\
% \texttt{last fret}    & \texttt{24} & Last fret shown (maximum 24).\\
% \texttt{width}        & \verb|\linewidth| & Total width of the diagram; any \TeX{} length.\\
% \texttt{accidentals}  & \texttt{both} & \texttt{sharp} (F$\sharp$), \texttt{flat} (G$\flat$) or \texttt{both} (stacked, F$\sharp$ over G$\flat$).\\
% \texttt{tuning}       & \texttt{D,B,G,D,G} & Five notes, string~1 to string~5. Write accidentals with \texttt{s}/\texttt{b}, e.g. \texttt{D,C,G,C,G}.\\
% \texttt{string order} & \texttt{tab} & \texttt{tab}: string~1 on top. \texttt{reverse}: string~5 on top.\\
% \texttt{left handed}  & false & Mirror the diagram (nut on the right).\\
% \texttt{string labels} & true & Show the string numbers 1--5.\\
% \texttt{fret numbers} & true & Show the fret numbers under the board.\\
% \texttt{markers}      & true & Show the inlay dots (frets 5, 7, 10, 12, 15, 17, 19, 22, 24).\\
% \texttt{focus}        & false & Grey out the unpainted notes when something is painted.\\
% \texttt{title}        & empty & Title printed above the diagram.\\
% \texttt{note font}    & \verb|\sffamily\bfseries| & Font commands for the note names.\\
% \bottomrule
% \end{tabular}
% \end{center}
%
% Boolean options can be written without a value (\texttt{focus} means
% \texttt{focus=true}).
%
% \BFexample{snip-flat}
% \BFexample{snip-lefty}
%
% \section{Tips and troubleshooting}
%
% \begin{itemize}
%   \item \textbf{Page size.}  24~frets with stacked note names are readable in
%     landscape (\texttt{pdflscape}, or a landscape \texttt{geometry}).  On
%     portrait pages use a window such as \texttt{last fret=12}, or
%     \texttt{accidentals=sharp}.
%   \item \textbf{Font size warnings.}  The note labels are scaled to the disc.
%     With the default Computer Modern fonts in OT1 encoding \LaTeX{} may
%     warn that some sizes are unavailable.  Loading \textsf{lmodern} (with
%     \textsf{fontenc}\,\texttt{T1}) or any scalable font family removes them.
%   \item \textbf{A \texttt{\#} inside another macro.}  In \verb|\bfpaintnote| a
%     literal \verb|#| works in the document body, but inside macro definitions
%     write \texttt{s} instead (\texttt{Fs}).
%   \item \textbf{Errors.}  Invalid positions, unknown note names and unknown
%     options give clear error messages and the offending item is ignored.
%   \item \textbf{Engines.}  Tested with pdf\LaTeX, Xe\LaTeX{} and Lua\LaTeX.
%   \item \textbf{Combining with text.}  The environment produces a centred block
%     and can be put in a \texttt{figure} float.
% \end{itemize}
%
% \section{Roadmap}
%
% Ideas for future versions: vertical (chord-box) orientation, fingering
% numbers inside the discs, scale-degree labels instead of note names, and
% ready-made scale and chord macros.  Contributions are welcome.
%
% \StopEventually{}
%
% \section{Implementation}
%
% This section documents the source code.  All internal functions and
% variables use the prefix \texttt{bfret} following the \textsf{expl3}
% conventions.  Diagrams are drawn at the end of the environment, so that the
% drawing order is guaranteed: board and frets, then strings, and finally the
% coloured discs with their note names.
%
%
%
% \subsection{Source code}
%
% \subsubsection{Package identification and dependencies}
%
%    \begin{macrocode}
%<*package>
%<package>\NeedsTeXFormat{LaTeX2e}[2022-06-01]
\RequirePackage{expl3}
\ProvidesExplPackage{banjofret}{2026/09/28}{1.0.0}
  {Five-string banjo fretboard diagrams with colour-coded chromatic notes}
\@ifl@t@r\fmtversion{2022-06-01}{}
  {\PackageError{banjofret}{LaTeX~2022-06-01~or~newer~is~required}{Please~update~your~TeX~distribution.}}
\ExplSyntaxOff
\RequirePackage{l3keys2e}
\RequirePackage{xcolor}
\RequirePackage{tikz}
\usetikzlibrary{calc}
\ExplSyntaxOn
%    \end{macrocode}
%
% \subsubsection{Messages}
%
%    \begin{macrocode}
\msg_new:nnn { banjofret } { unknown-key }
  { Unknown~option~'#1'. }
\msg_new:nnnn { banjofret } { outside-env }
  { #1~can~only~be~used~inside~a~'banjofret'~environment. }
  { Put~the~command~between~begin\{banjofret\}~and~end\{banjofret\}. }
\msg_new:nnnn { banjofret } { bad-note }
  { Unknown~note~name~'#1'. }
  { Use~a~letter~A--G~optionally~followed~by~\iow_char:N\# ~or~s~(sharp)~or~b~(flat),~e.g.~C,~Fs,~Bb. }
\msg_new:nnnn { banjofret } { bad-position }
  { Invalid~position~(#1,#2)~ignored. }
  { String~must~be~1--5,~fret~0--24;~string~5~only~exists~from~fret~5~onwards. }
\msg_new:nnnn { banjofret } { bad-pair }
  { Cannot~read~'#1'~as~(string,fret). }
  { Write~positions~as~(string,fret),~separated~by~semicolons,~e.g.~(3,5);(2,1). }
\msg_new:nnnn { banjofret } { bad-range }
  { Fret~window~must~satisfy~0<=first<=last<=24. }
  { Falling~back~to~the~full~fretboard~(0--24). }
\msg_new:nnnn { banjofret } { bad-tuning }
  { The~'tuning'~option~needs~exactly~five~notes~(strings~1--5). }
  { Falling~back~to~the~standard~open-G~tuning~D,B,G,D,G. }
\msg_new:nnn { banjofret } { legend-key-only }
  { Legend~entry~'#1'~has~no~label~and~was~ignored.~Use~colour=label. }
%    \end{macrocode}
%
% \subsubsection{Colour palette}
%
%    \begin{macrocode}
\definecolor{bfblue}{HTML}{7DB4E8}
\definecolor{bfgreen}{HTML}{86CFA3}
\definecolor{bfamber}{HTML}{F5C563}
\definecolor{bfcoral}{HTML}{F29E8E}
\definecolor{bfviolet}{HTML}{B9A3E3}
\definecolor{bfteal}{HTML}{7BCFCF}
\definecolor{bfgray}{HTML}{C9CED6}
\colorlet{bfroot}{bfcoral}
\colorlet{bfthird}{bfamber}
\colorlet{bffifth}{bfblue}
\colorlet{bfseventh}{bfviolet}
\colorlet{bfextension}{bfteal}
\colorlet{bftarget}{bfgreen}
\colorlet{bfpassing}{bfgray}
%    \end{macrocode}
%
% \subsubsection{Variables and constants}
%
%    \begin{macrocode}
\int_new:N \l__bfret_first_int
\int_new:N \l__bfret_last_int
\int_new:N \l__bfret_acc_int
\int_new:N \l__bfret_order_int
\int_new:N \l__bfret_dir_int
\int_new:N \l__bfret_xoff_int
\int_new:N \l__bfret_cols_int
\int_new:N \l__bfret_tmpa_int
\int_new:N \l__bfret_tmpc_int
\int_new:N \l__bfret_hash_int
\int_new:N \l__bfret_count_int
\bool_new:N \l__bfret_inenv_bool
\bool_new:N \l__bfret_open_bool
\bool_new:N \l__bfret_lefty_bool
\bool_new:N \l__bfret_slabels_bool
\bool_new:N \l__bfret_fnumbers_bool
\bool_new:N \l__bfret_markers_bool
\bool_new:N \l__bfret_focus_bool
\dim_new:N \l__bfret_w_dim
\dim_new:N \l__bfret_ss_dim
\dim_new:N \l__bfret_r_dim
\dim_new:N \l__bfret_d_dim
\dim_new:N \l__bfret_hh_dim
\dim_new:N \l__bfret_left_dim
\dim_new:N \l__bfret_right_dim
\dim_new:N \l__bfret_total_dim
\dim_new:N \l__bfret_margin_dim
\dim_new:N \l__bfret_fsa_dim
\dim_new:N \l__bfret_fsb_dim
\dim_new:N \l__bfret_xa_dim
\dim_new:N \l__bfret_ya_dim
\dim_new:N \l__bfret_xb_dim
\dim_new:N \l__bfret_yb_dim
\dim_new:N \l__bfret_tmpa_dim
\dim_new:N \l__bfret_tmpb_dim
\dim_new:N \l__bfret_ybot_dim
\tl_new:N \l__bfret_width_tl
\tl_new:N \l__bfret_title_tl
\tl_new:N \l__bfret_notefont_tl
\tl_new:N \l__bfret_tmpa_tl
\tl_new:N \l__bfret_key_tl
\tl_new:N \l__bfret_col_tl
\tl_new:N \l__bfret_draw_tl
\tl_new:N \l__bfret_txt_tl
\str_new:N \l__bfret_tmpa_str
\str_new:N \l__bfret_head_str
\str_new:N \l__bfret_orig_str
\str_new:N \l__bfret_tail_str
\clist_new:N \l__bfret_tuning_clist
\clist_new:N \l__bfret_names_clist
\seq_new:N \l__bfret_tuning_seq
\seq_new:N \l__bfret_legend_seq
\seq_new:N \l__bfret_tmpa_seq
\seq_new:N \l__bfret_tmpb_seq
\prop_new:N \l__bfret_paint_prop
\cs_new:Npn \__bfret_s: { \ensuremath { \sharp } }
\cs_new:Npn \__bfret_f: { \ensuremath { \flat } }
\seq_const_from_clist:Nn \c__bfret_sharp_seq
  { C , C\__bfret_s: , D , D\__bfret_s: , E , F , F\__bfret_s: , G , G\__bfret_s: , A , A\__bfret_s: , B }
\seq_const_from_clist:Nn \c__bfret_flat_seq
  { C , D\__bfret_f: , D , E\__bfret_f: , E , F , G\__bfret_f: , G , A\__bfret_f: , A , B\__bfret_f: , B }
\clist_const:Nn \c__bfret_natural_clist { 0 , 2 , 4 , 5 , 7 , 9 , 11 }
\clist_const:Nn \c__bfret_stringwidth_clist { 0.4pt , 0.55pt , 0.75pt , 1pt , 0.4pt }
\clist_const:Nn \c__bfret_markers_clist { 5 , 7 , 10 , 12 , 15 , 17 , 19 , 22 , 24 }
%    \end{macrocode}
%
% \subsubsection{Options}
%
%    \begin{macrocode}
\keys_define:nn { banjofret }
  {
    first~fret     .int_set:N  = \l__bfret_first_int ,
    first~fret     .initial:n  = 0 ,
    last~fret      .int_set:N  = \l__bfret_last_int ,
    last~fret      .initial:n  = 24 ,
    width          .tl_set:N   = \l__bfret_width_tl ,
    width          .initial:n  = \linewidth ,
    accidentals    .choices:nn = { sharp , flat , both }
      { \int_set:Nn \l__bfret_acc_int { \l_keys_choice_int } } ,
    accidentals    .initial:n  = both ,
    tuning         .clist_set:N = \l__bfret_tuning_clist ,
    tuning         .initial:n  = { D , B , G , D , G } ,
    string~order   .choices:nn = { tab , reverse }
      { \int_set:Nn \l__bfret_order_int { 3 - 2 * \l_keys_choice_int } } ,
    string~order   .initial:n  = tab ,
    left~handed    .bool_set:N = \l__bfret_lefty_bool ,
    left~handed    .initial:n  = false ,
    string~labels  .bool_set:N = \l__bfret_slabels_bool ,
    string~labels  .initial:n  = true ,
    fret~numbers   .bool_set:N = \l__bfret_fnumbers_bool ,
    fret~numbers   .initial:n  = true ,
    markers        .bool_set:N = \l__bfret_markers_bool ,
    markers        .initial:n  = true ,
    focus          .bool_set:N = \l__bfret_focus_bool ,
    focus          .initial:n  = false ,
    title          .tl_set:N   = \l__bfret_title_tl ,
    title          .initial:n  = { } ,
    note~font      .tl_set:N   = \l__bfret_notefont_tl ,
    note~font      .initial:n  = { \sffamily \bfseries } ,
    unknown        .code:n     =
      { \msg_error:nnx { banjofret } { unknown-key } { \l_keys_key_str } }
  }
\ProcessKeysOptions { banjofret }
\NewDocumentCommand \banjofretset { m }
  { \keys_set:nn { banjofret } {#1} }
%    \end{macrocode}
%
% \subsubsection{Note names and pitch classes}
%
%    \begin{macrocode}
\cs_new_protected:Npn \__bfret_note_to_pc:nN #1#2
  {
    \str_set:Nx \l__bfret_tmpa_str { \tl_to_str:n {#1} }
    \str_remove_all:Nn \l__bfret_tmpa_str { ~ }
    \str_set_eq:NN \l__bfret_orig_str \l__bfret_tmpa_str
    \str_set:Nx \l__bfret_head_str { \str_head:N \l__bfret_tmpa_str }
    \str_set:Nx \l__bfret_tail_str { \str_tail:N \l__bfret_tmpa_str }
    \exp_args:NV \str_case:nnF \l__bfret_head_str
      {
        { C } { \int_set:Nn #2 { 0 } }
        { D } { \int_set:Nn #2 { 2 } }
        { E } { \int_set:Nn #2 { 4 } }
        { F } { \int_set:Nn #2 { 5 } }
        { G } { \int_set:Nn #2 { 7 } }
        { A } { \int_set:Nn #2 { 9 } }
        { B } { \int_set:Nn #2 { 11 } }
      }
      {
        \msg_error:nnV { banjofret } { bad-note } \l__bfret_orig_str
        \int_set:Nn #2 { 0 }
      }
    \int_zero:N \l__bfret_hash_int
    \str_map_inline:Nn \l__bfret_tail_str
      {
        \str_case:nnF {##1}
          {
            { s } { \int_incr:N #2 }
            { b } { \int_decr:N #2 }
          }
          {
            \str_if_eq:nVTF {##1} \c_hash_str
              { \int_incr:N \l__bfret_hash_int }
              { \msg_error:nnV { banjofret } { bad-note } \l__bfret_orig_str }
          }
      }
    % A '#' character is doubled by \tl_to_str:n, hence the division by two.
    \int_add:Nn #2 { \int_div_truncate:nn { \l__bfret_hash_int } { 2 } }
    \int_set:Nn #2 { \int_mod:nn { \int_mod:nn {#2} { 12 } + 12 } { 12 } }
  }
\cs_new:Npn \__bfret_pc:nn #1#2
  {
    \int_mod:nn
      { \seq_item:Nn \l__bfret_tuning_seq {#1} + #2 - \int_compare:nTF { #1 = 5 } { 5 } { 0 } }
      { 12 }
  }
\cs_new_protected:Npn \__bfret_parse_tuning:
  {
    \seq_clear:N \l__bfret_tuning_seq
    \clist_map_inline:Nn \l__bfret_tuning_clist
      {
        \__bfret_note_to_pc:nN {##1} \l__bfret_tmpc_int
        \seq_put_right:NV \l__bfret_tuning_seq \l__bfret_tmpc_int
      }
    \int_compare:nF { \seq_count:N \l__bfret_tuning_seq = 5 }
      {
        \msg_error:nn { banjofret } { bad-tuning }
        \seq_set_split:Nnn \l__bfret_tuning_seq { , } { 2 , 11 , 7 , 2 , 7 }
      }
  }
%    \end{macrocode}
%
% \subsubsection{Colour resolution}
%
%    \begin{macrocode}
\cs_new:Npn \__bfret_color:n #1
  {
    \cs_if_exist:cTF { \c_backslash_str color@bf#1 } { bf#1 } {#1}
  }
%    \end{macrocode}
%
% \subsubsection{Painting commands}
%
%    \begin{macrocode}
\prg_new_conditional:Npnn \__bfret_valid:nn #1#2 { p , T , F , TF }
  {
    \bool_if:nTF
      {
        \int_compare_p:nNn {#1} > { 0 } && \int_compare_p:nNn {#1} < { 6 } &&
        \int_compare_p:nNn {#2} > { -1 } && \int_compare_p:nNn {#2} < { 25 } &&
        ( \int_compare_p:nNn {#1} < { 5 } || \int_compare_p:nNn {#2} > { 4 } )
      }
      { \prg_return_true: }
      { \prg_return_false: }
  }
\cs_new_protected:Npn \__bfret_check_env:n #1
  {
    \bool_if:NF \l__bfret_inenv_bool
      { \msg_error:nnn { banjofret } { outside-env } {#1} }
  }
\cs_new_protected:Npn \__bfret_put:nnn #1#2#3
  {
    \__bfret_valid:nnT {#1} {#2}
      {
        \tl_set:Nx \l__bfret_key_tl { \int_eval:n {#1} , \int_eval:n {#2} }
        \prop_put:NVn \l__bfret_paint_prop \l__bfret_key_tl {#3}
      }
  }
\cs_new_protected:Npn \__bfret_paint_pair:nn #1#2
  {
    \tl_if_blank:nF {#1}
      {
        \tl_set:Nn \l__bfret_tmpa_tl {#1}
        \tl_remove_all:Nn \l__bfret_tmpa_tl { ( }
        \tl_remove_all:Nn \l__bfret_tmpa_tl { ) }
        \seq_set_split:NnV \l__bfret_tmpb_seq { , } \l__bfret_tmpa_tl
        \int_compare:nTF { \seq_count:N \l__bfret_tmpb_seq = 2 }
          {
            \__bfret_valid:nnTF
              { \seq_item:Nn \l__bfret_tmpb_seq { 1 } }
              { \seq_item:Nn \l__bfret_tmpb_seq { 2 } }
              {
                \__bfret_put:nnn
                  { \seq_item:Nn \l__bfret_tmpb_seq { 1 } }
                  { \seq_item:Nn \l__bfret_tmpb_seq { 2 } }
                  {#2}
              }
              {
                \msg_error:nnxx { banjofret } { bad-position }
                  { \seq_item:Nn \l__bfret_tmpb_seq { 1 } }
                  { \seq_item:Nn \l__bfret_tmpb_seq { 2 } }
              }
          }
          { \msg_error:nnn { banjofret } { bad-pair } {#1} }
      }
  }
\NewDocumentCommand \bfpaint { m m }
  {
    \__bfret_check_env:n { \bfpaint }
    \seq_set_split:Nnn \l__bfret_tmpa_seq { ; } {#1}
    \seq_map_inline:Nn \l__bfret_tmpa_seq
      { \__bfret_paint_pair:nn {##1} {#2} }
  }
\cs_new_protected:Npn \__bfret_range:nNN #1#2#3
  {
    \seq_set_split:Nnn \l__bfret_tmpb_seq { - } {#1}
    \int_set:Nn #2 { \seq_item:Nn \l__bfret_tmpb_seq { 1 } }
    \int_compare:nTF { \seq_count:N \l__bfret_tmpb_seq > 1 }
      { \int_set:Nn #3 { \seq_item:Nn \l__bfret_tmpb_seq { 2 } } }
      { \int_set_eq:NN #3 #2 }
  }
\int_new:N \l__bfret_ra_int
\int_new:N \l__bfret_rb_int
\NewDocumentCommand \bfpaintstring { m m m }
  {
    \__bfret_check_env:n { \bfpaintstring }
    \__bfret_range:nNN {#2} \l__bfret_ra_int \l__bfret_rb_int
    \int_step_inline:nnn { \l__bfret_ra_int } { \l__bfret_rb_int }
      { \__bfret_put:nnn {#1} {##1} {#3} }
  }
\NewDocumentCommand \bfpaintfret { m m }
  {
    \__bfret_check_env:n { \bfpaintfret }
    \__bfret_range:nNN {#1} \l__bfret_ra_int \l__bfret_rb_int
    \int_step_inline:nnn { \l__bfret_ra_int } { \l__bfret_rb_int }
      {
        \int_step_inline:nn { 5 }
          { \__bfret_put:nnn { ####1 } { ##1 } {#2} }
      }
  }
\cs_new_protected:Npn \__bfret_paint_pc_string:nnn #1#2#3
  {
    \int_step_inline:nnn { 0 } { 24 }
      {
        \int_compare:nT { \__bfret_pc:nn {#3} { ##1 } = #1 }
          { \__bfret_put:nnn {#3} { ##1 } {#2} }
      }
  }
\NewDocumentCommand \bfpaintnote { m m }
  {
    \__bfret_check_env:n { \bfpaintnote }
    % Convert to a string first so that a literal '#' in the input is harmless.
    \clist_set:Nx \l__bfret_names_clist { \tl_to_str:n {#1} }
    \clist_map_inline:Nn \l__bfret_names_clist
      {
        \__bfret_note_to_pc:nN { ##1 } \l__bfret_tmpc_int
        \int_step_inline:nn { 5 }
          {
            \__bfret_paint_pc_string:nnn
              { \int_use:N \l__bfret_tmpc_int } {#2} { ####1 }
          }
      }
  }
\cs_new_protected:Npn \__bfret_legend_key:n #1
  { \msg_warning:nnn { banjofret } { legend-key-only } {#1} }
\cs_new_protected:Npn \__bfret_legend_pair:nn #1#2
  { \seq_put_right:Nn \l__bfret_legend_seq { { #1 } { #2 } } }
\NewDocumentCommand \bflegend { m }
  {
    \__bfret_check_env:n { \bflegend }
    \seq_clear:N \l__bfret_legend_seq
    \keyval_parse:NNn \__bfret_legend_key:n \__bfret_legend_pair:nn {#1}
  }
%    \end{macrocode}
%
% \subsubsection{Geometry helpers}
%
%    \begin{macrocode}
\cs_new:Npn \__bfret_xpos:n #1
  {
    \dim_eval:n
      {
        ( \int_eval:n { #1 - \l__bfret_xoff_int } \l__bfret_w_dim - 0.5 \l__bfret_w_dim )
        * \l__bfret_dir_int
      }
  }
\cs_new:Npn \__bfret_ypos:n #1
  { \dim_eval:n { \int_eval:n { \l__bfret_order_int * ( 3 - #1 ) } \l__bfret_ss_dim } }
\cs_new_protected:Npn \__bfret_seg:nnnnn #1#2#3#4#5
  {
    \dim_set:Nn \l__bfret_xa_dim { ( #1 ) * \l__bfret_dir_int }
    \dim_set:Nn \l__bfret_ya_dim { #2 }
    \dim_set:Nn \l__bfret_xb_dim { ( #3 ) * \l__bfret_dir_int }
    \dim_set:Nn \l__bfret_yb_dim { #4 }
    \draw [ #5 ]
      ( \l__bfret_xa_dim , \l__bfret_ya_dim ) --
      ( \l__bfret_xb_dim , \l__bfret_yb_dim ) ;
  }
\cs_new_protected:Npn \__bfret_fs:n #1
  { \exp_args:Nx \__bfret_fs_aux:n { \dim_to_decimal_in_unit:nn {#1} { 1pt } } }
\cs_new_protected:Npn \__bfret_fs_aux:n #1
  { \fontsize {#1} {#1} \selectfont }
\cs_new_protected:Npn \__bfret_setup_geometry:
  {
    \bool_if:nF
      {
        \int_compare_p:nNn { \l__bfret_first_int } > { -1 } &&
        \int_compare_p:nNn { \l__bfret_first_int } < { 25 } &&
        \int_compare_p:nNn { \l__bfret_last_int } > { \l__bfret_first_int - 1 } &&
        \int_compare_p:nNn { \l__bfret_last_int } < { 25 }
      }
      {
        \msg_error:nn { banjofret } { bad-range }
        \int_set:Nn \l__bfret_first_int { 0 }
        \int_set:Nn \l__bfret_last_int { 24 }
      }
    \int_set:Nn \l__bfret_xoff_int { \int_max:nn { \l__bfret_first_int - 1 } { 0 } }
    \bool_set:Nn \l__bfret_open_bool
      { \int_compare_p:nNn { \l__bfret_first_int } = { 0 } }
    \int_set:Nn \l__bfret_cols_int
      { \l__bfret_last_int - \l__bfret_xoff_int + \bool_if:NTF \l__bfret_open_bool { 1 } { 0 } }
    \int_set:Nn \l__bfret_dir_int { \bool_if:NTF \l__bfret_lefty_bool { -1 } { 1 } }
    \dim_set:Nn \l__bfret_total_dim { \l__bfret_width_tl }
    \dim_set:Nn \l__bfret_margin_dim
      { \bool_if:NTF \l__bfret_slabels_bool { 14pt } { 0pt } }
    \dim_set:Nn \l__bfret_w_dim
      { ( \l__bfret_total_dim - \l__bfret_margin_dim ) / \l__bfret_cols_int }
    \dim_set:Nn \l__bfret_ss_dim
      { \dim_min:nn { 0.8 \l__bfret_w_dim } { 26pt } }
    \dim_set:Nn \l__bfret_r_dim
      { \dim_min:nn { \l__bfret_w_dim } { \l__bfret_ss_dim } }
    \dim_set:Nn \l__bfret_r_dim { 0.44 \l__bfret_r_dim }
    \dim_set:Nn \l__bfret_d_dim { 2 \l__bfret_r_dim }
    \dim_set:Nn \l__bfret_fsa_dim { 1.15 \l__bfret_r_dim }
    \dim_set:Nn \l__bfret_fsb_dim { 0.8 \l__bfret_r_dim }
    \dim_set:Nn \l__bfret_hh_dim { 2.6 \l__bfret_ss_dim }
    \dim_set:Nn \l__bfret_left_dim
      { \bool_if:NTF \l__bfret_open_bool { - \l__bfret_w_dim } { 0pt } }
    \dim_set:Nn \l__bfret_right_dim
      { \int_eval:n { \l__bfret_last_int - \l__bfret_xoff_int } \l__bfret_w_dim }
  }
%    \end{macrocode}
%
% \subsubsection{Drawing the fretboard}
%
%    \begin{macrocode}
\cs_new_protected:Npn \__bfret_draw_fretlines:
  {
    \int_step_inline:nnn { \l__bfret_xoff_int } { \l__bfret_last_int }
      {
        \dim_set:Nn \l__bfret_tmpa_dim
          { \int_eval:n { ##1 - \l__bfret_xoff_int } \l__bfret_w_dim }
        \bool_if:nTF
          { \int_compare_p:nNn {##1} = { \l__bfret_xoff_int } && \bool_if_p:N \l__bfret_open_bool }
          {
            \__bfret_seg:nnnnn
              { \l__bfret_tmpa_dim } { - \l__bfret_hh_dim }
              { \l__bfret_tmpa_dim } { \l__bfret_hh_dim }
              { line~width = 3pt , line~cap = butt }
          }
          {
            \__bfret_seg:nnnnn
              { \l__bfret_tmpa_dim } { - \l__bfret_hh_dim }
              { \l__bfret_tmpa_dim } { \l__bfret_hh_dim }
              { line~width = 0.6pt }
          }
      }
  }
\cs_new_protected:Npn \__bfret_draw_strings:
  {
    \int_step_inline:nn { 5 }
      {
        \bool_if:nT
          { \int_compare_p:nNn {##1} < { 5 } || \int_compare_p:nNn { \l__bfret_last_int } > { 4 } }
          {
            \dim_set:Nn \l__bfret_tmpb_dim
              {
                \bool_if:nTF
                  { \int_compare_p:nNn {##1} = { 5 } && \int_compare_p:nNn { \l__bfret_first_int } < { 6 } }
                  { \int_eval:n { 5 - \l__bfret_xoff_int } \l__bfret_w_dim - 0.5 \l__bfret_w_dim }
                  { \l__bfret_left_dim }
              }
            \__bfret_seg:nnnnn
              { \l__bfret_tmpb_dim } { \__bfret_ypos:n {##1} }
              { \l__bfret_right_dim } { \__bfret_ypos:n {##1} }
              { line~width = \clist_item:Nn \c__bfret_stringwidth_clist {##1} }
          }
      }
  }
\cs_new_protected:Npn \__bfret_label:
  {
    \tl_use:N \l__bfret_notefont_tl
    \clist_if_in:NVTF \c__bfret_natural_clist \l__bfret_tmpc_int
      {
        \__bfret_fs:n { \l__bfret_fsa_dim }
        \seq_item:Nn \c__bfret_sharp_seq { \l__bfret_tmpc_int + 1 }
      }
      {
        \int_case:nn { \l__bfret_acc_int }
          {
            { 1 }
              {
                \__bfret_fs:n { \l__bfret_fsa_dim }
                \seq_item:Nn \c__bfret_sharp_seq { \l__bfret_tmpc_int + 1 }
              }
            { 2 }
              {
                \__bfret_fs:n { \l__bfret_fsa_dim }
                \seq_item:Nn \c__bfret_flat_seq { \l__bfret_tmpc_int + 1 }
              }
            { 3 }
              {
                \__bfret_fs:n { \l__bfret_fsb_dim }
                \shortstack [ c ]
                  {
                    \seq_item:Nn \c__bfret_sharp_seq { \l__bfret_tmpc_int + 1 } \\
                    \seq_item:Nn \c__bfret_flat_seq { \l__bfret_tmpc_int + 1 }
                  }
              }
          }
      }
  }
\cs_new_protected:Npn \__bfret_disc:nn #1#2
  {
    \__bfret_valid:nnT {#1} {#2}
      {
        \dim_set:Nn \l__bfret_xa_dim { \__bfret_xpos:n {#2} }
        \dim_set:Nn \l__bfret_ya_dim { \__bfret_ypos:n {#1} }
        \int_set:Nn \l__bfret_tmpc_int { \__bfret_pc:nn {#1} {#2} }
        \prop_get:NnNTF \l__bfret_paint_prop { #1 , #2 } \l__bfret_tmpa_tl
          {
            \tl_set:Nx \l__bfret_col_tl { \__bfret_color:n { \l__bfret_tmpa_tl } }
            \node
              [ circle , draw = black , fill = \l__bfret_col_tl , line~width = 0.6pt ,
                inner~sep = 0pt , minimum~size = \l__bfret_d_dim , text = black ]
              at ( \l__bfret_xa_dim , \l__bfret_ya_dim )
              { \__bfret_label: } ;
          }
          {
            \bool_lazy_and:nnTF
              { \bool_if_p:N \l__bfret_focus_bool }
              { ! \prop_if_empty_p:N \l__bfret_paint_prop }
              {
                \node
                  [ circle , draw = black!30 , fill = white , line~width = 0.5pt ,
                    inner~sep = 0pt , minimum~size = \l__bfret_d_dim , text = black!40 ]
                  at ( \l__bfret_xa_dim , \l__bfret_ya_dim )
                  { \__bfret_label: } ;
              }
              {
                \node
                  [ circle , draw = black , fill = white , line~width = 0.5pt ,
                    inner~sep = 0pt , minimum~size = \l__bfret_d_dim , text = black ]
                  at ( \l__bfret_xa_dim , \l__bfret_ya_dim )
                  { \__bfret_label: } ;
              }
          }
      }
  }
\cs_new_protected:Npn \__bfret_draw_discs:
  {
    \int_step_inline:nn { 5 }
      {
        \int_step_inline:nnn { \l__bfret_first_int } { \l__bfret_last_int }
          { \__bfret_disc:nn { ##1 } { ####1 } }
      }
  }
\cs_new_protected:Npn \__bfret_draw_labels:
  {
    \bool_if:NT \l__bfret_slabels_bool
      {
        \int_step_inline:nn { 5 }
          {
            \dim_set:Nn \l__bfret_xa_dim
              { ( \l__bfret_left_dim - 4pt ) * \l__bfret_dir_int }
            \dim_set:Nn \l__bfret_ya_dim { \__bfret_ypos:n {##1} }
            \node
              [ anchor = \bool_if:NTF \l__bfret_lefty_bool { west } { east } ,
                inner~sep = 0pt , text = black ]
              at ( \l__bfret_xa_dim , \l__bfret_ya_dim )
              { \sffamily \footnotesize ##1 } ;
          }
      }
  }
\cs_new_protected:Npn \__bfret_draw_frettrack:
  {
    \dim_set:Nn \l__bfret_ybot_dim { - \l__bfret_hh_dim }
    \bool_if:NT \l__bfret_fnumbers_bool
      {
        \dim_sub:Nn \l__bfret_ybot_dim { 9pt }
        \int_step_inline:nnn { \int_max:nn { \l__bfret_first_int } { 1 } } { \l__bfret_last_int }
          {
            \dim_set:Nn \l__bfret_xa_dim { \__bfret_xpos:n {##1} }
            \node [ inner~sep = 0pt , text = black ]
              at ( \l__bfret_xa_dim , \l__bfret_ybot_dim )
              { \sffamily \footnotesize ##1 } ;
          }
      }
    \bool_if:NT \l__bfret_markers_bool
      {
        \dim_sub:Nn \l__bfret_ybot_dim { 9pt }
        \clist_map_inline:Nn \c__bfret_markers_clist
          {
            \bool_if:nT
              {
                \int_compare_p:nNn {##1} > { \l__bfret_first_int - 1 } &&
                \int_compare_p:nNn {##1} < { \l__bfret_last_int + 1 }
              }
              {
                \dim_set:Nn \l__bfret_xa_dim { \__bfret_xpos:n {##1} }
                \int_compare:nTF { \int_mod:nn {##1} { 12 } = 0 }
                  {
                    \dim_set:Nn \l__bfret_xb_dim { \l__bfret_xa_dim - 3pt }
                    \fill ( \l__bfret_xb_dim , \l__bfret_ybot_dim ) circle ( 1.8pt ) ;
                    \dim_set:Nn \l__bfret_xb_dim { \l__bfret_xa_dim + 3pt }
                    \fill ( \l__bfret_xb_dim , \l__bfret_ybot_dim ) circle ( 1.8pt ) ;
                  }
                  { \fill ( \l__bfret_xa_dim , \l__bfret_ybot_dim ) circle ( 1.8pt ) ; }
              }
          }
      }
  }
\cs_new_protected:Npn \__bfret_legend_item:nn #1#2
  {
    \tl_set:Nx \l__bfret_col_tl { \__bfret_color:n {#1} }
    \node
      [ circle , draw = black , fill = \l__bfret_col_tl , line~width = 0.6pt ,
        inner~sep = 0pt , minimum~size = 10pt , anchor = west ]
      ( bfret-lc ) at ( bfret-legpos ) { } ;
    \node
      [ anchor = west , inner~xsep = 4pt , inner~ysep = 0pt , text = black ]
      ( bfret-lt ) at ( bfret-lc.east ) { \sffamily \footnotesize #2 } ;
    \coordinate ( bfret-legpos ) at ( $ ( bfret-lt.east ) + ( 12pt , 0pt ) $ ) ;
  }
\cs_new_protected:Npn \__bfret_draw_legend:
  {
    \seq_if_empty:NF \l__bfret_legend_seq
      {
        \dim_sub:Nn \l__bfret_ybot_dim { 14pt }
        \dim_set:Nn \l__bfret_xa_dim
          { \bool_if:NTF \l__bfret_lefty_bool { - \l__bfret_right_dim } { \l__bfret_left_dim } }
        \coordinate ( bfret-legpos ) at ( \l__bfret_xa_dim , \l__bfret_ybot_dim ) ;
        \seq_map_inline:Nn \l__bfret_legend_seq { \__bfret_legend_item:nn ##1 }
      }
  }
\cs_new_protected:Npn \__bfret_draw_title:
  {
    \tl_if_blank:VF \l__bfret_title_tl
      {
        \dim_set:Nn \l__bfret_xa_dim
          { ( \l__bfret_left_dim + \l__bfret_right_dim ) / 2 * \l__bfret_dir_int }
        \dim_set:Nn \l__bfret_ya_dim { \l__bfret_hh_dim + 6pt }
        \node [ anchor = south , inner~sep = 0pt , text = black ]
          at ( \l__bfret_xa_dim , \l__bfret_ya_dim )
          { \sffamily \bfseries \large \l__bfret_title_tl } ;
      }
  }
%    \end{macrocode}
%
% \subsubsection{The banjofret environment}
%
%    \begin{macrocode}
\cs_new_protected:Npn \__bfret_begin:n #1
  {
    \prop_clear:N \l__bfret_paint_prop
    \seq_clear:N \l__bfret_legend_seq
    \bool_set_true:N \l__bfret_inenv_bool
    \keys_set:nn { banjofret } {#1}
    \__bfret_parse_tuning:
    \__bfret_setup_geometry:
  }
\cs_new_protected:Npn \__bfret_end:
  {
    \par \medskip \noindent \hfil
    \begin{tikzpicture}
      \__bfret_draw_title:
      \__bfret_seg:nnnnn
        { \l__bfret_left_dim } { \l__bfret_hh_dim }
        { \l__bfret_right_dim } { \l__bfret_hh_dim } { line~width = 0.8pt }
      \__bfret_seg:nnnnn
        { \l__bfret_left_dim } { - \l__bfret_hh_dim }
        { \l__bfret_right_dim } { - \l__bfret_hh_dim } { line~width = 0.8pt }
      \__bfret_draw_fretlines:
      \__bfret_draw_strings:
      \__bfret_draw_discs:
      \__bfret_draw_labels:
      \__bfret_draw_frettrack:
      \__bfret_draw_legend:
    \end{tikzpicture}
    \hfil \par \medskip
  }
\NewDocumentEnvironment { banjofret } { O{} }
  { \__bfret_begin:n {#1} }
  { \__bfret_end: }
%</package>
%    \end{macrocode}
%
% \Finale
\endinput
